% ---------------------------------------------------------------------------
%  IV. Methodology  — the paper's primary contribution
%  Floats:  figures/fig-vantage.tex   (Fig. 3, vantage points and blind spots)
%           tables/tab-phases.tex     (Table II, seven-phase pipeline)
%  Rule:    this section states the METHOD only. Hypothesis verdicts and any
%           measured outcome belong in section 06. Keep it that way: a
%           reviewer noticed the mixture in the previous version.
% ---------------------------------------------------------------------------

\section{Methodology}
\label{sec:method}

\subsection{Principles}

Four rules govern the procedure.

\begin{itemize}[leftmargin=1.2em]
\item \textbf{P1: observe before publishing.} We listen passively before injecting anything,
because there is no way to predict how an undocumented safety-critical system will react.
\item \textbf{P2: verify effects, not acknowledgments.} A command counts as validated only
once an independent state topic changes. This is the working form of H4.
\item \textbf{P3: triangulate.} We accept a finding only when two independent sources agree,
drawn from network traces, introspection, decompiled application code and field logs.
\item \textbf{P4: remain non-destructive.} No firmware modification, no root, no persistent
change to vendor partitions. This is an ethical commitment, and with a single unit it is also
what makes the work repeatable.
\end{itemize}

\subsection{Vantage points and their blind spots}

The whole methodology exists because no single observation point tells us enough
(Fig.~\ref{fig:vantage}). Introspection lists names, but says nothing about preconditions.
Traffic capture shows what the vendor application happened to send during one session, which is
not the same as what the chassis would accept. Decompiled code gives intent and the implicit
setup steps, but leaves open whether a given path is reachable on this firmware. Triangulating
between the three is not a stylistic preference; it is the only way to cover what any one of
them misses.

% Fig. 3 — vantage points and their blind spots. Used by sections/04-methodology.
% This figure carries the methodological argument: three partial views, one
% ground truth. Uses the shared styles and the \figsub macro from preamble.tex.
\begin{figure}[!t]
\centering
\resizebox{0.94\columnwidth}{!}{%
\begin{tikzpicture}[x=1cm, y=1cm]

  \def\colA{0} \def\colB{3.35} \def\colC{6.7}

  \node[figbox blue, minimum width=2.5cm] (a) at (\colA,0)
    {Introspection\\[1pt]\figsub{names, types}};
  \node[figbox teal, minimum width=2.5cm] (b) at (\colB,0)
    {Traffic capture\\[1pt]\figsub{framing, order}};
  \node[figbox rust, minimum width=2.5cm] (c) at (\colC,0)
    {Application code\\[1pt]\figsub{intent, setup}};

  \node[font=\scriptsize\itshape, text=figslate, align=center] (a2) at (\colA,-1.25)
    {blind to\\preconditions};
  \node[font=\scriptsize\itshape, text=figslate, align=center] (b2) at (\colB,-1.25)
    {blind to\\unused commands};
  \node[font=\scriptsize\itshape, text=figslate, align=center] (c2) at (\colC,-1.25)
    {blind to\\reachability};

  \draw[figarrow] (a) -- (a2);
  \draw[figarrow] (b) -- (b2);
  \draw[figarrow] (c) -- (c2);

  \node[figbox slate, minimum width=8.6cm, minimum height=0.7cm] (t) at (\colB,-2.5)
    {Field telemetry: the only ground truth for execution (P2, H4)};

  \draw[figarrow] (a2) -- (t);
  \draw[figarrow] (b2) -- (t);
  \draw[figarrow] (c2) -- (t);

\end{tikzpicture}%
}
\caption{Each vantage point answers a different question and is blind to the others. A
recovered command is accepted only after state telemetry confirms a physical effect.}
\label{fig:vantage}
\end{figure}


\subsection{The seven-phase pipeline}

Table~\ref{tab:phases} gives, for each phase, the evidence it produces, the hypothesis it bears
on, and the condition under which we stopped. The stopping criteria are not a formality. With
several hundred topics on the wire, an exploration without them simply does not converge.

% Table II — the seven-phase pipeline. Used by sections/04-methodology.
% Rule: every phase must declare the evidence it produces AND the condition
% under which it terminates. A phase without a stopping criterion is a task,
% not a method step, and reviewers said so about the previous version.
\begin{table}[!t]
\caption{Reverse-engineering pipeline}
\label{tab:phases}
\centering
\scriptsize
\setlength{\tabcolsep}{3pt}
\begin{tabular}{@{}p{0.15cm}p{1.25cm}p{2.35cm}p{0.5cm}p{2.15cm}@{}}
\toprule
\textbf{\#} & \textbf{Phase} & \textbf{Evidence produced} & \textbf{Hyp.} & \textbf{Stopping criterion} \\
\midrule
1 & Network mapping & Segment map, open ports & --- & All hosts and reachable ports enumerated \\
2 & Introspection & Topic and service inventory & H1 & Inventory stable across sessions \\
3 & Passive broker & Payload classes on the broker & H2 & No command payload after repeated sessions \\
4 & Application audit & Names, types, implicit setup steps & H1 & Every observed name located in code \\
5 & Frame capture & Message framing and ordering & H1 & Client reproduces vendor framing \\
6 & Edge integration & On-device services and bridges & --- & Capability reachable without a host PC \\
7 & Field validation & State transitions, pose deltas & H3, H4 & Effect confirmed on independent topic \\
\bottomrule
\end{tabular}
\end{table}


We kept the negative results instead of discarding them: a legacy velocity topic with no
subscriber, two candidate speech topics, three HTTP ports on the head, and the broker command
path of H2. Each one closed a branch of the search, which is worth as much as an open one.

We report the cost of the method because a procedure nobody can budget for is of little use to
anyone else. Phases~1--6 took two weeks and phase~7 a further week, spread over eight sessions
with the physical robot and carried out by one full-time engineer. We decompiled two
applications, and archived full introspections rather than counts, which turns the phase~2
stopping criterion into something testable.

Within a session the inventory is exactly reproducible: two introspections \SI{45}{\second} apart
returned set-identical results, \num{448} topics, \num{334} services and \num{64} nodes. Across
sessions it is not, but it varies in a structured way. An inventory taken three weeks earlier is
a strict subset of the later one: nothing disappeared, and the \num{44} topics and \num{13}
services that appeared belong without exception to one subsystem, the downward-facing depth
camera, which was not running on the earlier date. The upward-facing camera's \num{115} topics
were identical on both dates. The phase~2 criterion should therefore be read as stability under a
fixed set of running subsystems, not as a single number; an integrator who counts once and
assumes the graph is closed will miss whatever was idle that day. This also explains a still
earlier count of \num{455} topics, taken before we archived inventories and which we can no
longer diff, but which falls in the range the mechanism predicts.

The exploration scripts, the recovered interface description and the
logs behind every measured result reported here are released with this
paper~\cite{cybel2026}; the two indicators in Table~\ref{tab:metrics} that carry no archived
log are marked as such rather than implying an archive that does not exist.
